\section{ENERGIA E QUANTITÀ DI MOTO DEL CAMPO ELETTROMAGNETICO}

\subsection{L'energia del campo elettromagnetico}

Nel caso elettrostatico l'energia del campo è stata definita come il lavoro che occorre fare sulle cariche per portarle nella configurazione finale: siccome alla fine le cariche sono ferme questo lavoro, e quindi l'energia del campo, è l'energia potenziale del sistema. Per esempio, se carichiamo un condensatore con una pila ed una resistenza, il lavoro è quello (positivo) fatto dalla pila più quello (negativo) fatto, sempre sugli elettroni, dalle forze ``viscose'' esercitate dal reticolo entro cui gli elettroni si muovono.\\
Nel caso generale, in cui nella configurazione finale le cariche sono in moto, il lavoro che occorre fare per realizzare la configurazione comprende necessariamente anche quello per fornire alle cariche e ai loro supporti (cilindri conduttori rotanti, ioni in movimento in una soluzione, \dots) la necessaria energia: nel caso dei cilindri rotanti o di spire in movimento questa è solo energia cinetica, ma più in generale (si pensi ad una resistenza che si scalda) si tratta di energia interna.\\
A parità di cariche e di correnti l'energia totale del sistema dipende quindi anche da caratteristiche non elettriche del sistema, come il materiale (rame, alluminio, \dots) di cui sono fatti i cilindri rotanti o la massa degli ioni, o la temperatura delle resistenze: dato che non possiamo prescindere dal compiere almeno parte di questo lavoro ``meccanico'', dobbiamo includerlo nella definizione dell'energia del campo? Abbiamo visto come la dinamica relativistica, più precisamente il 3$^\circ$ principio, dia sostanza al concetto di campo (particolarmente nei casi non statici), imponendo di considerare il campo come un sistema per conto suo: il nostro sistema di cariche e di correnti interagenti tramite il campo elettromagnetico deve quindi essere visto come formato da due sistemi: uno è il sistema meccanico costituito dalle cariche (pensate come particelle che non interagiscono direttamente fra di loro) e dai loro supporti, e l'altro è il campo elettromagnetico (con cui le cariche interagiscono). Di conseguenza, anche l'energia totale è la somma delle energie dei singoli sistemi, cioè l'energia \underline{cinetica} (o l'\underline{energia interna}) del sistema meccanico più l'energia del campo, che risulta così definita come la differenza fra l'energia totale del ``sistema composto'', e l'energia del ``sistema meccanico''.\\
In sostanza, a prima vista non sembra che ci sia molta differenza rispetto all'affermazione che l'energia del sistema solare è la somma delle energie cinetiche dei pianeti e dell'energia potenziale del sistema: la differenza fondamentale sta nel fatto che mentre l'energia potenziale di un sistema newtoniano di particelle ad ogni istante dipende dalle posizioni (eventualmente anche dalle velocità) delle particelle a quello stesso istante (e ciò consente di considerare la somma dell'energia cinetica e dell'energia potenziale come l'energia di un unico sistema, il sistema di particelle interagenti), l'energia del campo elettromagnetico, a causa della velocità finita di propagazione delle interazioni, dipende dalle posizioni e velocità delle cariche a tutti gli istanti precedenti e quindi non può essere espressa in termini della configurazione istantanea del sistema meccanico, mentre invece, come vedremo fra poco, dipende dai campi solo all'istante considerato: è questo fatto che suggerisce (o impone) di considerare particelle (non interagenti) e campo
elettromagnetico come due sistemi distinti.\\
Questa discussione presuppone che sia possibile una chiara distinzione fra campo e sistema meccanico, per cui dovremo riprenderla nei capitoli 14 e 15 dove tratteremo i campi elettrico e magnetico nella materia: in questi casi la materia viene considerata come un sistema termodinamico, per cui si perde traccia della sua struttura interna in termini di cariche (elettroni, protoni) e dei campi a queste associate. Quindi il concetto di energia del campo elettromagnetico al quale arriviamo in questo capitolo in effetti si applica, in teoria a qualsiasi sistema purché considerato da un punto di vista microscopico (non termodinamico), in pratica a sistemi di poche cariche o anche a sistemi la cui struttura interna non viene modificata apprezzabilmente dai campi (p.es.\ i conduttori): questo è quello che normalmente viene chiamato l'\textbf{\textit{elettromagnetismo nel vuoto}}, per distinguerlo da quello nella materia.

In definitiva, quella che (nell'elettromagnetismo nel vuoto) definiamo \textbf{energia del campo elettromagnetico} corrispondente ad una certa configurazione di cariche e di correnti, in generale non è tutto il lavoro $L^\mathrm{est}$ fatto dall'esterno sulle sorgenti per realizzare la configurazione, ma $L^\mathrm{est}$ meno la variazione $\Delta E_m$ dell'energia (cinetica o interna) del sistema meccanico:
\begin{equation}
U\defeq L^\mathrm{est}-\Delta E_m
\tag{11.1}
\end{equation}
per cui, se con $L^\mathrm{em}$ indichiamo il lavoro fatto sul sistema meccanico dalle forze elettromagnetiche, cioè dal campo elettrico (la forza di Lorentz è a lavoro nullo), grazie al teorema delle forze vive (o al teorema delle forze vive generalizzato)
\begin{equation}
L^\mathrm{est}+L^\mathrm{em}=\Delta E_m\quad\ \Rightarrow\quad\ U=-L^\mathrm{em}
\tag{11.2}
\end{equation}
cioè l'energia del campo è definita come l'opposto del lavoro fatto sul sistema meccanico dalle forze elettromagnetiche (quelle che nella fisica newtoniana sono le ``forze interne'' al sistema di particelle interagenti), o, come si suol dire, il lavoro fatto dalle forze esterne contro le forze del campo, per realizzare il campo stesso.\\
È bene tuttavia avere presente che se $\Delta E_m$ è solo la variazione dell'energia cinetica degli elettroni in un conduttore, questa è del tutto trascurabile rispetto all'energia del campo (faremo un confronto numerico nell'esempio che discuteremo più avanti), per cui in questi casi l'energia del campo praticamente coincide con il lavoro fatto dall'esterno sul sistema.

Il lavoro fatto nell'unità di tempo (potenza) dalle forze del campo sulle cariche, se $\vec u_i$ è il campo di velocità delle cariche di densità $\rho_i$, siccome il campo magnetico esercita forze a lavoro nullo, è dato da
\begin{equation}
W(t)=\textstyle\sum_i\displaystyle\int\rho_i\vec E\cdot\vec u_i\,dV=\int\vec E(x,y,z,t)\cdot\vec J(x,y,z,t)\,dV
\tag{11.3}
\end{equation}
e l'energia del campo è quindi
\begin{equation}
U(t)=-\int_{-\infty}^tW(t')\,dt'.
\tag{11.4}
\end{equation}
Vogliamo esprimere l'energia in termini dei campi (come abbiamo fatto nel caso elettrostatico): moltiplichiamo scalarmente la 4$^\mathrm{a}$ e la 2$^\mathrm{a}$ delle (6.53) rispettivamente per $\vec E$
e per $\vec B$ e sottraiamole:
\begin{equation}
\vec E\cdot\rot\vec B-\vec B\cdot\rot\vec E=\frac{4\pi}{c}\vec E\cdot\vec J+\frac{1}{c}\Big(\vec E\cdot\frac{\partial\vec E}{\partial t}+\vec B\cdot\frac{\partial\vec B}{\partial t}\Big)\cdot
\tag{11.5}
\end{equation}
Si ha l'identità
\begin{equation}
\vec a\cdot\rot\vec b-\vec b\cdot\rot\vec a=-\dive(\vec a\wedge\vec b\,)
\tag{11.6}
\end{equation}
che si dimostra o mediante calcolo diretto, oppure \big(stesse notazioni come nella (7.19) e (7.20)\big)
\begin{equation}
\begin{aligned}
\dive(\vec a\wedge\vec b\,)&=\epsilon_{ijk}\partial_i(a_jb_k)=\epsilon_{ijk}\big[(\partial_ia_j)b_k+(\partial_ib_k)a_j\big]\\
&=\vec b\cdot\rot\vec a-\vec a\cdot\rot\vec b
\end{aligned}
\tag{11.7}
\end{equation}
(abbiamo utilizzato l'antisimmetria del simbolo di Levi Civita).\\
Grazie alla (11.7) dalla (11.5) si ha
\begin{equation}
-\vec E\cdot\vec J=\frac{1}{8\pi}\frac{\partial}{\partial t}\Big(\vec E^2+\vec B^2\Big)+\frac{c}{4\pi}\dive\big(\vec E\wedge\vec B\,\big).
\tag{11.8}
\end{equation}
Integriamo la (11.8) su tutto lo spazio; per il teorema di Green l'ultimo termine eguaglia il flusso all'infinito di $\vec E\wedge\vec B$: i campi (se le sorgenti sono al finito) si annullano abbastanza rapidamente all'infinito, per cui detto flusso si annulla, quindi per le (11.3) e (11.4)
\begin{equation}
U(t)=\frac{1}{8\pi}\int\big(\vec E^2+\vec B^2\big)\,dV
\tag{11.9}
\end{equation}
che è l'espressione cercata per l'energia del campo; essa si riduce alla (1.94) nel caso elettrostatico, e dimostra la (8.43) nel caso magnetostatico. La (11.9) mostra che l'energia del campo elettromagnetico è una \underline{funzione di stato}, cioè che la definizione di energia del campo elettromagnetico è una ``buona definizione''.\\
Equivalentemente, dalle (11.3) e (11.4),
\begin{equation}
-\frac{dU(t)}{dt}=\int\vec E\cdot\vec J\,dV
\tag{11.10}
\end{equation}
cioè la diminuzione dell'energia nell'unità di tempo eguaglia il lavoro fatto nell'unità di tempo dalle forze del campo sulle cariche; se il sistema sorgenti (compresi i loro supporti meccanici)--campo è un sistema isolato (e quindi le forze del campo sono le sole che agiscono sulle sorgenti), per il teorema delle forze vive (generalizzato) il lavoro fatto nell'unità di tempo sulle cariche è la derivata dell'energia $E_m(t)$ del sistema meccanico, e la (11.10) esprime la conservazione dell'energia per il sistema (isolato) costituito dalle cariche e dal campo:
\begin{equation}
\frac{d}{dt}\Big(U(t)+E_m(t)\Big)=0\,.
\tag{11.11}
\end{equation}
In generale, cioè se il sistema non è isolato, dalla (11.1) si ha
\begin{equation}
\frac{d}{dt}\Big(U(t)+E_m(t)\Big)=W^\mathrm{est}.
\tag{11.12}
\end{equation}

\paragraph{Esempi}
\begin{enumerate}
\item Consideriamo un circuito, come quello della figura 8.4, costituito da un generatore di f.e.m.\ $V_0$, un'induttanza $L$ e una resistenza $R$. Se prendiamo come sistema meccanico solo le cariche di conduzione, le forze esterne sono quelle (non conservative) del generatore (per esempio la forza di gravità in una pila a gravità) che fanno lavoro positivo sulle cariche, e le forze esercitate dal reticolo sugli elettroni (``forze viscose'') che fanno lavoro negativo: le prime sviluppano una potenza \ $W_G(t)=V_0I(t)$, e le seconde \ $-W_\mathrm{J}(t)=-RI^2(t)$. Stimiamo l'energia cinetica degli elettroni: supponiamo che il circuito sia costituito da 1\,m di filo conduttore di sezione 1\,mm, (volume complessivo 1\,cm$^3$), percorso alla fine da una corrente di 1\,A. In queste condizioni \big(v.\ (3.5)\big) la velocità degli elettroni è di circa $10^{-2}$\,cm/s, quindi la loro energia cinetica totale è
\begin{equation}
E_c=\frac{1}{2}mu^2\cdot n\approx10^{-27}\cdot10^{-4}\cdot10^{23}=10^{-8}\,\mathrm{erg}
\tag{11.13}
\end{equation}
e vedremo fra poco che questa è trascurabile rispetto all'energia del campo, per cui tutto il lavoro delle forze esterne va ad incrementare l'energia del campo. Quindi
\begin{equation}
\frac{dU}{dt}=W_G-W_\mathrm{J}=V_0I(t)-RI^2(t)\,.
\tag{11.14}
\end{equation}
Siccome per un circuito percorso da una corrente lentamente variabile
\begin{equation}
\int\vec E\cdot\vec J\,dV=\mathcal E(t)I(t)=-L\dot I(t)I(t)
\tag{11.15}
\end{equation}
dalla (11.10) si ottiene
\begin{equation}
U(t)=\frac{1}{2}LI^2(t)
\tag{11.16}
\end{equation}
in accordo con la (8.42); inoltre, dalle (11.14) e (11.15) ritroviamo l'equazione del circuito (8.55).\\
Raggiunte le condizioni stazionarie ($t\to\infty$, $I$ costante), l'energia del campo non cambia più e si ha quindi $W_G=W_\mathrm{J}$.\\
Per confrontare l'energia del campo con la (11.13), supponiamo (è il caso peggiore) che il circuito non contenga ``deliberatamente'' un'induttanza (cioè bobine, solenoidi), e quindi che $L$ sia solo dovuta al fatto che qualunque circuito ha necessariamente un coefficiente di autoinduzione: in queste condizioni prendiamo per $L$ l'espressione (8.32) (coefficiente di autoinduzione di un toro), senza il fattore $N^2$, quindi (meglio un cattivo numero che nessun numero) $L\approx10^{-8}$\,H, per cui dalla (11.16)
\begin{equation}
U\approx10^{-8}\,\mathrm{J}\approx10^7\,E_c.
\tag{11.17}
\end{equation}
Ovviamente avremmo potuto prendere come sistema meccanico l'intero circuito, comprendente quindi anche il reticolo in cui si muovono gli elettroni; in questo caso le forze esterne, oltre a quelle del generatore, sono (a livello microscopico) quelle che determinano gli scambi di calore (Joule) con l'ambiente esterno: il bilancio energetico espresso dalla (11.14) non viene modificato. Se il circuito è
termicamente isolato, allora \ $W_\mathrm{J}\equiv RI^2$ \ nella (11.14) è la variazione nell'unità di tempo dell'energia interna del circuito.

\item Riprendiamo in considerazione il problema trattato nel capitolo 7.6 della forza esercitata dal campo magnetico di un solenoide finito (e fisso) su una spira percorsa da corrente, e supponiamo che tanto il solenoide quanto la spira siano alimentati da generatori $G_1$ e $G_2$ che mantengono costante la corrente rispettivamente nel solenoide (1) e nella spira (2). Le forze esterne al sistema spira+solenoide--campo sono quelle esercitate dai generatori. Abbiamo visto che quando la spira si muove tutto il lavoro fatto dal generatore $G_2$ oltre a quello che viene dissipato per effetto Joule eguaglia la variazione dell'energia cinetica della spira. Inoltre, a causa del moto della spira, il solenoide è soggetto ad un campo magnetico (quello della spira) variabile nel tempo, quindi in esso viene indotta una f.e.m.\ $\mathcal E_1(t)$ (negativa se la spira si avvicina la solenoide, cioè se le correnti sono concordi, positiva altrimenti), ed il generatore $G_1$ deve quindi sviluppare (in più) una potenza \ $W_1=-\mathcal E_1(t)I_1$ \ per mantenere costante la corrente nel solenoide. Questo lavoro va a modificare l'energia del campo magnetico: infatti l'energia magnetica è data da
\begin{equation}
U_M=\frac{1}{8\pi}\int\big(\vec B_1+\vec B_2\big)^2\,dV=\frac{1}{8\pi}\int\big(\vec B_1^2+\vec B_2^2+2\vec B_1\cdot\vec B_2\big)\,dV
\tag{11.18}
\end{equation}
e se cambia la posizione relativa della spira e del solenoide, cambia l'ultimo termine nella (11.18), e quindi l'energia del campo. La cosa può essere vista in termini quantitativi nel modo seguente. Per la (8.42)
\begin{equation}
U_M=\frac{1}{2}L_{11}I_1^2+\frac{1}{2}L_{22}I_2^2+L_{12}I_1I_2
\tag{11.19}
\end{equation}
e siccome $L_{11}$, $L_{22}$, $I_1$ e $I_2$ sono costanti (ma non lo è il coefficiente di mutua induzione $L_{12}$),
\begin{equation}
\frac{dU_M(t)}{dt}=\frac{dL_{12}}{dt}I_1I_2=\frac{dL_{21}I_2}{dt}I_1=\frac{1}{c}\frac{d\Phi_{21}(\vec B_2)}{dt}I_1=-\mathcal E_1(t)I_1=W_1.
\tag{11.20}
\end{equation}
In questo caso, quindi, le forze esterne modificano sia l'energia cinetica del sistema meccanico, sia l'energia del campo: la potenza $W_1+W_2$ sviluppata dai due generatori ($W_1=W_2$ grazie al teorema di reciprocità) è uguale alla derivata rispetto al tempo della somma dell'energia del campo e dell'energia cinetica della spira.

\item La forza che si esercita fra due cariche ferme uguali è repulsiva, e l'energia del campo (elettrico) aumenta se le due cariche vengono avvicinate: i due fatti sono ovviamente conseguenza l'uno dell'altro.\\
Se consideriamo due spire coassiali uguali ed entrambe percorse nello stesso verso dalla corrente $I$, è ancora vero che l'energia del campo (magnetico) aumenta al diminuire della distanza fra le spire, in quanto il coefficiente di mutua induzione cresce avvicinando le spire: se le due spire sono a distanza infinita l'energia del campo è il doppio dell'energia della singola spira (il coefficiente di mutua induzione è nullo), mentre se le due spire sono sovrapposte, l'energia del campo è quella di una spira percorsa dalla corrente $2I$ e quindi è quattro volte quella della spira singola (il coefficiente di mutua induzione è uguale a quello di autoinduzione). Però, e
ciò a prima vista può apparire paradossale, a differenza del caso delle due cariche, la forza fra le due spire è attrattiva.\\
Non c'è alcun paradosso: se la corrente nelle due spire è costante, allora in ogni spira deve essere inserito un generatore di corrente e, come nell'esempio precedente, sia l'aumento dell'energia del campo, sia l'aumento dell'energia cinetica delle due spire (se sono libere di avvicinarsi) sono dovuti al lavoro fatto dai generatori.\\
Se invece nelle due spire, che supponiamo superconduttrici, non sono presenti generatori, allora a causa della mutua induzione la corrente nelle spire diminuisce in misura tale che, nonostante le due spire si avvicinino, anche l'energia del campo diminuisce e corrispondentemente aumenta la loro energia cinetica.\\
Analizziamo questo secondo caso (spire superconduttrici) calcolando l'energia cinetica $E_c$ delle due spire quando sono sovrapposte, se inizialmente si trovano a distanza infinita.\\
Detto $L$ il coefficiente di autoinduzione di ogni spira e $M$ quello di mutua induzione (che dipende dalla distanza fra le spire), $I_1(t)$ e $I_2(t)$ le correnti nelle spire, siccome le due spire più il campo costituiscono un sistema isolato, nell'ipotesi che il moto delle spire sia lento (approssimazione delle correnti lentamente variabili) dalla (11.11) si ha
\begin{equation}
\frac{d}{dt}\Big(\frac{1}{2}LI_1^2+\frac{1}{2}LI_2^2+MI_1I_2+E_c\Big)=0\,.
\tag{11.21}
\end{equation}
Le equazioni dei due circuiti sono
\begin{equation}
L\dot I_1+\frac{d\,MI_2}{dt}=0;\qquad L\dot I_2+\frac{d\,MI_1}{dt}=0\,.
\tag{11.22}
\end{equation}
Dalle (11.22), se $I_0$ è la corrente che scorre inizialmente in ciascuna delle due spire ed essendo all'inizio $M=0$, si ha
\begin{equation}
\begin{aligned}
&LI_1+MI_2=LI_0;\qquad LI_2+MI_1=LI_0\quad\ \Rightarrow\\
&I_1(t)=I_2(t)\equiv I(t)=\frac{L}{L+M}I_0\,.
\end{aligned}
\tag{11.23}
\end{equation}
Poiché alla fine $M=L$, dalla (11.21) e dalla (11.23) si ha
\begin{equation}
E_c=\frac{1}{2}LI_0^2+\frac{1}{2}LI_0^2-\frac{1}{2}\cdot4L\Big(\frac{I_0}{2}\Big)^2=\frac{1}{2}LI_0^2
\tag{11.24}
\end{equation}
quindi la corrente in ogni spira si è dimezzata, di conseguenza si è dimezzata l'energia del campo (alla fine si ha una sola spira con corrente $I_0$), e metà dell'energia iniziale si è trasformata in energia cinetica delle spire.\\
Le (11.22) esprimono il fatto che la f.e.m.\ in una spira superconduttrice è nulla, e quindi il flusso del campo magnetico attraverso la spira resta costante (``conservazione del flusso'', come espresso dalle prime due equazioni nella (11.23)); questo fatto permette di arrivare in modo più elegante (anche se completamente equivalente a quello che abbiamo seguito sopra) alla (11.24), in quanto se il flusso attraverso ciascuna spira quando queste sono sovrapposte deve essere uguale al flusso iniziale, la corrente totale nelle due spire deve eguagliare la corrente iniziale $I_0$ in ciascuna spira.
\end{enumerate}

\subsection{Il teorema di Poynting}

La (11.8), che esprime in forma locale il bilancio energetico, è nota come \textbf{teorema di Poynting}, del quale vogliamo ora discutere il significato.\\
Definiamo rispettivamente la \textbf{densità di energia} del campo ed il \textbf{vettore di Poynting} come
\begin{equation}
u(x,y,z,t)\defeq\frac{1}{8\pi}\big(\vec E^2+\vec B^2\big);\qquad\quad\vec S(x,y,z,t)\defeq\frac{c}{4\pi}\vec E\wedge\vec B
\tag{11.25}
\end{equation}
e quindi la (11.8) si scrive
\begin{equation}
\dive\vec S+\frac{\partial u}{\partial t}=-\vec E\cdot\vec J.
\tag{11.26}
\end{equation}
Dividiamo lo spazio in due parti $V_1$ e $V_2$. Posto
\begin{equation}
U_{V_1}\equiv\int_{V_1}u(x,y,z,t)\,dV;\qquad U_{V_2}\equiv\int_{V_2}u(x,y,z,t)\,dV
\tag{11.27}
\end{equation}
dalla (11.26) si ha
\begin{equation}
-\Big(\frac{dU_{V_1}}{dt}+\int_{V_1}\vec E\cdot\vec J\Big)=\Phi_{V_1}(\vec S\,);\qquad\frac{dU_{V_2}}{dt}+\int_{V_2}\vec E\cdot\vec J=-\Phi_{V_2}(\vec S\,).
\tag{11.28}
\end{equation}
Supponiamo (solo per semplicità) che non ci siano forze esterne agenti sul sistema, allora per il teorema delle forze vive \big(vedi (11.11)\big) la (11.28) si può riscrivere come
\begin{equation}
-\frac{d}{dt}\Big(U_{V_1}+E_{V_1}^c\Big)=\Phi_{V_1}(\vec S\,);\qquad\frac{d}{dt}\Big(U_{V_2}+E_{V_2}^c\Big)=-\Phi_{V_2}(\vec S\,)
\tag{11.29}
\end{equation}
dove \ $E_{V_1}^c$, $E_{V_2}^c$ \ sono rispettivamente le energie cinetiche della parte del sistema (meccanico) contenuta nel volume $V_1$ e di quella contenuta nel volume $V_2$. Tenuto conto che
\begin{equation}
\Phi_{V_2}(\vec S\,)=-\Phi_{V_1}(\vec S\,)
\tag{11.30}
\end{equation}
la (11.29) costituisce la forma globale del teorema di Poynting e si legge nel modo seguente:

\textit{la diminuzione nell'unità di tempo dell'energia totale contenuta dentro un volume $V$ è data dal flusso uscente da $V$ del vettore di Poynting; questo flusso coincide con il flusso entrante nel volume circostante ed eguaglia l'aumento dell'energia totale in esso contenuta.}

Il flusso del vettore di Poynting è quindi l'energia trasferita fra due sistemi nell'unità di tempo, quello contenuto nel volume $V_1$ e quello contenuto nel volume adiacente $V_2$, ed il vettore di Poynting ha quindi il significato di \textbf{densità di corrente di energia}, così come $\vec J$ ha il significato di densità di corrente di carica. L'analogia fra $\rho,\vec J$ e $u,\vec S$ è molto stretta: se all'interno del volume $V$ non sono presenti sorgenti, la (11.29) dice che l'energia del campo contenuta in $V$ può cambiare solo perché c'è un flusso di energia attraverso la superficie che delimita $V$: all'interno del volume $V$ vale l'equazione di continuità per l'energia
\begin{equation}
\dive\vec S+\frac{\partial u}{\partial t}=0\qquad\quad(x,y,z\in V)
\tag{11.31}
\end{equation}
identica all'equazione (3.10) di continuità per la corrente.\\
L'introduzione del concetto di densità e flusso di energia del campo elettromagnetico non ha solo un valore formale: permette di considerare il campo (ed eventualmente le sorgenti) contenuto in un volume $V$ esso stesso come un sistema in interazione con il campo e le sorgenti esterni a $V$, ed il vettore di Poynting descrive questa interazione: esso permette di andare oltre l'idea maxwelliana di localizzazione dell'energia, in quanto ora non solo sappiamo quanta energia si trova in un dato volume ma anche, seguendo le linee di flusso del vettore $\vec S$, dove essa andrà o da dove essa è venuta; per la (11.26) pozzi e sorgenti di (flusso di) energia sono i punti in cui \ $\partial u/\partial t+\vec E\cdot\vec J\ne0$.

\paragraph{Esempio}
\begin{quote}
Consideriamo un circuito costituito da una pila e da un filo conduttore di resistività $\varrho_c$. Prendiamo come volume $V$ quello (strettamente) occupato dal circuito, pila compresa. All'interno del conduttore è presente un campo elettrico che per la legge di Ohm in forma locale è dato dalla (3.17); se con $E_t$ indichiamo la componente di $\vec E$ parallela al filo (l'unica non nulla dentro al conduttore), si ha
\begin{equation}
E_t=\varrho_cJ=\frac{R_{\Delta L}}{\Delta L}I
\tag{11.32}
\end{equation}
dove $R_{\Delta L}$ è la resistenza di un tratto di filo di lunghezza (infinitesima) $\Delta L$.\\
Siccome la componente tangenziale di $\vec E$ \ è continua \ ($\rot\vec E=0$) attraverso la superficie del filo, la (11.32) è anche la componente tangenziale di $\vec E$ sulla superficie esterna di $V$. Poiché il filo è percorso da corrente, è presente anche un campo magnetico che, sempre sulla superficie esterna di $V$, vale
\begin{equation}
B_\theta=\frac{2I}{ac}
\tag{11.33}
\end{equation}
dove $a$ è il raggio del conduttore. Allora in ogni punto della superficie esterna di $V$ il vettore di Poynting è non nullo, ha una componente diretta verso il filo (componente radiale), che per la (11.32) e la (11.33) vale
\begin{equation}
S_\varrho=-\frac{c}{4\pi}E_tB_\theta=-\frac{R_{\Delta L}I^2}{2\pi a\Delta L}
\tag{11.34}
\end{equation}
e nel tratto di filo di lunghezza $\Delta L$ si ha un flusso \underline{entrante} di energia dato da
\begin{equation}
\Phi(\vec S\,)=2\pi a\,\Delta L\,S_\varrho=R_{\Delta L}I^2
\tag{11.35}
\end{equation}
e quindi nell'intero filo entra la potenza \ $W_\mathrm{J}\equiv RI^2$, esattamente uguale a quella che viene dissipata per effetto Joule. Ovviamente non è un caso: siccome dentro il filo non cambia né l'energia cinetica degli elettroni, né il campo, l'energia che entra o si trasforma in energia interna del filo se questo è termicamente isolato, oppure viene ritrasferita sotto forma di calore all'esterno (avremo quindi anche un flusso di energia termica uscente). Inoltre, siccome l'energia del campo all'esterno del volume $V$ (cioè del circuito) non cambia, a questo flusso entrante in $V$ deve corrispondere un uguale flusso uscente: infatti all'interno della pila la corrente è opposta al campo elettrico, quindi campo elettrico e campo magnetico danno luogo ad un flusso del vettore di Poynting \underline{uscente} dalla pila, che quindi deve
essere uguale a quello entrante nel filo. In definitiva, le linee di forza del vettore di Poynting escono dalla pila ed entrano nel conduttore: la pila è la sorgente ed il filo il pozzo.\\
Se poi nel circuito è presente anche un'induttanza, come prima troviamo che nel filo entra la potenza \ $W_\mathrm{J}\equiv RI^2$ \ ma, a differenza di prima, dalla pila deve uscire un flusso maggiore, dato che l'energia del campo magnetico deve aumentare: infatti, se facciamo il confronto con il caso precedente a parità di corrente nel filo, e quindi dentro la pila, e quindi a parità di campo magnetico ovunque e di campo elettrico nel filo, il campo elettrico dentro la pila è maggiore in quanto riceve anche il contributo del campo elettrico autoindotto, concorde con il campo elettrico (conservativo) presente all'interno della pila.
\end{quote}

Ovviamente gli esempi che abbiamo appena discusso non aggiungono nulla di nuovo a ciò che già sapevamo a proposito del bilancio energetico, anche perché si tratta di situazioni stazionarie o quasi--stazionarie; forniscono però un modo diverso di vedere le cose, un modo in cui il campo gioca un ruolo importante e che diventa insostituibile nei casi in cui interviene in maniera essenziale la propagazione dei campi (campi rapidamente variabili).

\subsection{La quantità di moto del campo elettromagnetico}

Per costruire il campo elettromagnetico non solo occorre fare lavoro sulle sorgenti ma, in generale, occorre anche cedere dall'esterno ad esse un impulso; come per l'energia, parte di questo impulso si ritrova come quantità di moto del sistema meccanico (sorgenti e loro supporti), ed il resto è, per definizione, la \textbf{quantità di moto del campo elettromagnetico}. La quantità di moto del sistema meccanico è la somma dell'impulso totale ceduto dalle forze esterne e di quello $\vec{\mathcal I}^\mathrm{em}$ ceduto dalle forze del campo (``forze interne''): quest'ultimo sarebbe nullo se valesse il 3$^\circ$ principio della dinamica nella sua formulazione newtoniana (il risultante delle forze interne -- in questo caso elettromagnetiche -- fra i componenti del sistema meccanico sarebbe nullo), ma in generale è diverso da zero ed il suo opposto è quindi la quantità di moto $\vec P$ del campo. Quindi
\begin{equation}
\vec P(t)=-\vec{\mathcal I}^\mathrm{em}(-\infty,t)=-\int_{-\infty}^tdt'\int\big(\rho\vec E+\frac{1}{c}\vec J\wedge\vec B\,\big)\,dV\,.
\tag{11.36}
\end{equation}
L'integrale di volume che compare nella (11.36) si trasforma in una espressione contenente solo i campi sostituendo alle sorgenti $\rho,\vec J$ le loro espressioni ricavate dalle equazioni di Maxwell: il calcolo è alquanto più laborioso che non quello che ha portato alla (11.9) (anche in questo caso i flussi all'infinito sono nulli), per cui lo omettiamo; il risultato è
\begin{equation}
-\int\big(\rho\vec E+\frac{1}{c}\vec J\wedge\vec B\,\big)\,dV=\frac{d}{dt}\Big(\frac{1}{4\pi c}\int\vec E\wedge\vec B\,dV\Big)=\frac{1}{c^2}\frac{d}{dt}\int\vec S\,dV
\tag{11.37}
\end{equation}
e quindi
\begin{equation}
\vec P=\frac{1}{4\pi c}\int\vec E\wedge\vec B\,dV=\frac{1}{c^2}\int\vec S\,dV\,.
\tag{11.38}
\end{equation}
La quantità di moto del campo è quindi proporzionale all'integrale del vettore di Poynting: \ $\vec S/c^2$ \ è perciò la \textbf{densità di quantità di moto} del campo.
Se sul sistema sorgenti--campo non agiscono forze esterne, cioè se il sistema è isolato, la quantità di moto totale si conserva; infatti
\begin{equation}
\begin{aligned}
&\frac{d\vec P^\mathrm{\,mecc}}{dt}=\int\big(\rho\vec E+\frac{1}{c}\vec J\wedge\vec B\,\big)\,dV=-\frac{d\vec P}{dt}\quad\ \Rightarrow\\
&\frac{d}{dt}\big(\vec P^\mathrm{\,mecc}+\vec P\,\big)=0\,.
\end{aligned}
\tag{11.39}
\end{equation}

\subsection{Un esempio \dots\ mancato, e il problema dei 4/3}

Il problema che vogliamo discutere in questo paragrafo, anche se verrà trattato in forma elementare, ci porta un po' oltre a quello che è l'obiettivo di queste lezioni: infatti riguarda una problematica che, pur avendo radici nel passato, è ancora oggetto in tempi recenti di discussione nella letteratura specializzata.\\
Consideriamo il sistema di figura 11.1: fra le armature di un condensatore a facce piane parallele, carico, è presente un solenoide, con l'asse parallelo alle facce del condensatore, percorso da corrente costante. Il tutto è montato su un carrello.
\fig{fig11_1}{11.1}
Nel sistema è presente sia un campo elettrico che un campo magnetico che, a parte gli effetti di bordo (che possono essere resi trascurabili), sono fra loro ortogonali. Quindi il vettore di Poynting non è nullo ed è ortogonale sia all'asse del solenoide che a quello del condensatore, per cui il campo elettromagnetico possiede una quantità di moto.\\
Supponiamo ora che il condensatore o il solenoide venga ``scaricato'' su una resistenza, anch'essa parte di tutto il dispositivo. Alla fine il campo elettrico o quello magnetico è nullo, per cui è nullo il vettore di Poynting e quindi la quantità di moto del campo. Dove è andata a finire la quantità di moto che il campo inizialmente possedeva? La conclusione più naturale, visto che la quantità di moto del campo più quella del sistema meccanico (condensatore+solenoide+carrello) si conserva, è che si sia trasferita al sistema meccanico, e quindi che il carrello si metta in moto fino a raggiungere una certa velocità, che poi mantiene. Questa conclusione sembra corroborata dalla considerazione seguente: supponiamo che venga ``scaricato'' il solenoide; siccome la corrente varia, il campo elettrico indotto dalla variazione del campo magnetico esercita una forza sulle cariche presenti su ciascuna armatura del condensatore: le componenti di queste forze parallele alle due armature sono uguali, in quanto le cariche sono opposte, ma anche la componente di $\vec E$ parallela alle armature è opposta su di esse. Quindi il condensatore, solidale al carrello, è soggetto ad una forza parallela alle armature e ortogonale all'asse del solenoide (come il vettore di Poynting). Per contro, a prima vista, il campo elettrico del condensatore non esercita alcuna forza sul solenoide, e quindi il risultante delle forze (interne!) fra il condensatore ed il solenoide è diverso da zero, ed il carrello si dovrebbe mettere in moto. Il discorso è analogo, e la conclusione la stessa, se invece viene scaricato il condensatore: in questo caso è il campo magnetico prodotto dalla corrente di spostamento che esercita una forza sulla corrente che scorre nel solenoide.

Prima di andare avanti, mettiamo qualche numero per renderci conto degli ordini di
grandezza: supponiamo di avere in un volume $V$ di $10^4$\,cm$^3$ un campo elettrico di $3\times10^3$\,volt/cm ed un campo magnetico di $10^3$\,G, fra loro ortogonali. La quantità di moto del campo vale
\begin{equation}
P=\frac{1}{4\pi c}\,E\,B\,V\approx\frac{10\cdot10^3\cdot10^4}{4\pi\cdot3\times10^{10}}\approx3\times10^{-4}\,\mathrm{g\,cm/s}
\tag{11.40}
\end{equation}
e quindi anche se tutto il dispositivo avesse una massa di soli 10\,g, raggiungerebbe la velocità di $3\times10^{-5}$\,cm/s.

È evidente che non possiamo fare un esperimento per vedere se il carrello si muove o no, e quindi per decidere se il meccanismo che abbiamo ipotizzato è corretto: per fortuna esiste un bellissimo e semplicissimo ragionamento, dovuto a P.~Menotti, che ci permette di escludere che il carrello possa mettersi in moto. Supponiamo il carrello fermo e almeno uno dei due campi (p.es.\ quello magnetico) nullo; se il meccanismo sopra ipotizzato fosse corretto, ``caricando'' il solenoide (cioè facendo circolare corrente in esso) il sistema dovrebbe mettersi in moto, e successivamente ``scaricando'' il solenoide il sistema dovrebbe rallentare fino a fermarsi (il sistema è isolato e la quantità di moto all'inizio e alla fine è nulla), quindi in definitiva il sistema si sarebbe spostato di un certo tratto, in una certa direzione. Se osserviamo tutto il processo da un riferimento in moto con velocità $\vec v$ rispetto al laboratorio in direzione ortogonale allo spostamento del carrello, ci accorgiamo che all'inizio e alla fine il momento angolare del sistema (rispetto ad un polo arbitrario) è diverso: la velocità del sistema è sempre $-\vec v$, ma la distanza dal polo è cambiata. Ma il momento angolare di un sistema isolato si conserva (in ogni riferimento), quindi il carrello non si è mosso.

Si noti che questo ragionamento fa uso soltanto delle leggi di conservazione della quantità di moto e del momento angolare, e del principio di relatività, e sembra portare ad una conclusione assurda: se il carrello non si muove, da dove viene la quantità di moto del campo elettromagnetico quando nel solenoide scorre la corrente (e il condensatore è carico)? Sembra proprio (ed è) un rompicapo: l'unica conclusione logica che salvi sia le leggi di conservazione che il principio di relatività (e per la quale indicheremo anche possibili meccanismi) è che, così come un campo elettromagnetico \underline{statico} può possedere una quantità di moto, anche un sistema meccanico \underline{fermo} possa possedere una quantità di moto (chiamata nella letteratura ``hidden momentum'', cioè ``quantità di moto nascosta'', in quanto non associata ad un moto del centro di massa del sistema): in effetti un sistema meccanico (nel nostro caso il condensatore e il solenoide) per non disintegrarsi richiede la presenza di forze di coesione, quindi in esso è presente un campo di forze (del quale non è necessario specificare la natura, anche se sostanzialmente è di natura elettromagnetica), e quindi non è assurdo che in seguito alle tensioni interne al materiale provocate dall'azione del campo elettromagnetico sulle sorgenti anche a questo campo possa essere associata una quantità di moto. Quindi in definitiva la quantità di moto totale del sistema è nulla sia all'inizio che durante tutto il processo: la quantità di moto del campo elettromagnetico è compensata, istante per istante, dalla quantità di moto ``nascosta'' del sistema meccanico, e perciò non era corretta la nostra supposizione che durante la ``scarica'' (o la ``carica'') del solenoide si esercitasse solo una forza sulle armature del condensatore, e non una reazione sul solenoide.

Un altro possibile meccanismo per cui in un sistema meccanico in cui è presente sia un campo elettrico che una corrente (e quindi un campo magnetico) si può instaurare una quantità di moto nascosta è ben illustrato dal seguente modello.\\
Immaginiamo di avere un tubo in cui circolano liberamente in versi opposti cariche positive e negative e che questa ``spira'', percorsa da una corrente stazionaria, si trovi in un campo elettrico uniforme, come nella figura 11.2. Supponiamo che la corrente giri in senso antiorario \ $(A\to C\to B\to D\to A)$.
\fig{fig11_2}{11.2}
Fissiamo l'attenzione sulle cariche positive: per effetto del campo elettrico la loro velocità nella metà di destra è maggiore che nella metà di sinistra, tuttavia siccome la corrente è stazionaria la loro densità è inversamente proporzionale alla velocità ($\rho v$ è costante), per cui, se prendiamo per la quantità di moto di una particella l'espressione non relativistica $\vec p=m\vec v$, si vede subito che la quantità di moto totale $\vec P$ delle cariche è nulla. Ma se prendiamo l'espressione relativistica \ $\vec p=m\gamma(v)\vec v$ \ ne risulta una quantità di moto totale diretta verso l'alto nella figura 11.2 perché, grazie al fatto che $\gamma(v)$ aumenta con la velocità, i punti in cui la velocità è maggiore portano un contributo maggiore alla quantità di moto (rispetto al caso newtoniano), nonostante la minore densità (che è sempre inversamente proporzionale a $v$). Le cariche negative portano un contributo alla quantità di moto identico a quello delle cariche positive. Questa è la ``quantità di moto nascosta'' nel sistema, in quanto tutto il sistema è stazionario (il centro di massa è fermo), e il calcolo esplicito mostrerebbe che è esattamente uguale e opposto alla quantità di moto del campo elettromagnetico: siccome per la (11.26) le sorgenti del vettore di Poynting $\vec S$ sono i punti in cui \ $\vec E\cdot\vec J<0$ \ (nel nostro caso i campi sono statici per cui $\partial u/\partial t=0$), le linee di forza di $\vec S$ vanno dalla parte superiore $(B-D-A)$ della spira (dove $\vec E$ e $\vec J$ formano un angolo ottuso) a quella inferiore $(A-C-B)$ e quindi il verso della quantità di moto del campo è effettivamente opposto a quello delle cariche circolanti.

Come abbiamo detto all'inizio del paragrafo, questa problematica nasce da lontano: all'inizio di questo secolo (1900), praticamente in concomitanza con la nascita della teoria della relatività, Abraham e Lorentz tentarono di costruire una teoria puramente elettromagnetica dell'elettrone, attribuendone la massa, in base alla formula $E=mc^2$, esclusivamente all'energia elettrostatica. Questa teoria si scontrò subito con alcune difficoltà: in primo luogo siccome l'energia è inversamente proporzionale al raggio $a$ dell'elettrone (pensato come una sferetta), non è possibile considerare l'elettrone come un oggetto elementare puntiforme senza struttura (cioè prendere il limite $a\to0$), ma è necessario attribuire ad esso una dimensione dell'ordine di \ $e^2/mc^2\approx2.8\times10^{-13}$\,cm (il cosiddetto ``raggio classico dell'elettrone''). Ma se l'elettrone ha una struttura, cioè delle dimensioni, la speranza di poterlo descrivere in termini puramente elettromagnetici deve essere abbandonata, in quanto necessariamente devono essere presenti forze di natura non elettromagnetica che gli impediscano di disintegrarsi per effetto della repulsione coulombiana.\\
Ma il problema principale nacque da un evidente conflitto con la teoria della relatività: se schematizziamo l'elettrone come una sferetta di raggio $a$ con carica uniformemente
distribuita sulla superficie (ma le difficoltà sono indipendenti dal modello), l'energia e la quantità di moto del campo elettromagnetico di un elettrone in moto con velocità $\vec v$ risultano essere dati da
\begin{equation}
E=\gamma(v)\frac{e^2}{2a}\,\big(1+\frac{1}{3}\beta^2\big);\qquad\vec p=\frac{4}{3}\gamma(v)\frac{e^2}{2a}\,\vec v
\tag{11.41}
\end{equation}
e quindi, se la massa dell'elettrone è di natura puramente elettromagnetica,
\begin{equation}
\begin{aligned}
&m=E(v=0)/c^2=\frac{e^2}{2ac^2}\quad\ \Rightarrow\\
&E=m\gamma c^2\big(1+\frac{1}{3}\beta^2\big),\qquad\vec p=\frac{4}{3}m\gamma(v)\vec v
\end{aligned}
\tag{11.42}
\end{equation}
ed è evidente che energia e quantità di moto non si trasformano correttamente come un quadrivettore, e in particolare non vale la relazione (5.21) \ $m^2c^4=E^2-p^2c^2$.

Alla soluzione di questo problema, noto come ``il problema dei 4/3'', hanno contribuito in molti, da Poincaré a Fermi (quando ancora lavorava a Pisa): le forze di coesione di natura non elettromagnetica necessariamente presenti all'interno dell'elettrone per garantirne la stabilità costituiscono un campo (il \textbf{\textit{campo degli sforzi di Poincaré}}) al quale è associata energia e quantità di moto (nascosta), e solo quando anche questi contributi vengono presi in considerazione si ottengono per l'elettrone una energia e quantità di moto totali (elettromagnetici + meccanici) che si trasformano correttamente come un quadrivettore.\\
La conclusione è quindi che la struttura meccanica di un sistema ed il campo di forze che ne garantisce la stabilità è sempre importante, in particolare quando si vogliano costruire modelli classici delle particelle.
